\chapterpage{Introduction}
\mychapter{Introduction}

\section{Background}

Public transportation is one of the essential components of today's urban fabric. Buses are the fastest growing mode of transport and carry the greatest proportion of daily commuters globally, especially in the developing world, which is undergoing rapid urbanisation \cite{itf_outlook_2023}. It is crucial that passenger occupancy information is accurate in order to optimize scheduling, fleet management, prevent overcrowding, and ensure safety for passengers. Typical APC systems use hardware sensors including an infrared beam, a pressure mat, or a ticketing system. These methods are common, but they can also have drawbacks such as high installation and maintenance expenses, susceptibility to environmental conditions, and decreased accuracy with multiple passengers simultaneously entering or exiting the system. Field evaluations comparing infrared, WiFi, and video-based APC solutions confirm that traditional sensing technologies show accuracy declines under crowded or high-throughput boarding conditions~\cite{grgurevic2021review}. With recent progress in computer vision and deep learning, camera-based passenger counting systems have been developed. Many buses are already using surveillance cameras for security reasons, making vision an affordable and scalable solution \cite{Chris2023}. Neural networks can accurately assess occupancy in complex visual scenes, and deep learning models that integrate object detection and density estimation have presented state-of-the-art results in crowd counting benchmarks \cite{gao2022caaisurvey}.

\section{Problem Statement}

The ability to count passengers with high accuracy is a key need in intelligent public transportation systems. Accurate passenger data allows transport authorities and bus companies to make informed decisions, optimize routes, ensure passenger safety, and monitor bus occupancy. Traditional passenger counting mechanisms such as infrared sensors, pressure sensors, and ticketing systems tend to be inaccurate as a result of malfunctions, simultaneous boarding and alighting, and passenger overlap \cite{Aleksander2024}. In recent years, computer vision and deep learning have shown promising results in the field of crowd counting and object detection related to surveillance tasks \cite{sindagi2022survey}. In the case of bus passenger counting, however, it is not so easy to apply these methods due to the special conditions found in the interior of a bus. The passengers are tightly packed, so occlusion is a serious issue, in which people partially or even completely cover each other \cite{sundararaman2021}. In addition, surveillance cameras are usually installed on the ceiling, resulting in perspective distortion and making passenger heads appear small, especially in the rear sections of the bus \cite{rapid2020}. Illumination, movement of vehicles, and frequent passenger movements add to the difficulty in accurately counting \cite{Claire2021}. Existing crowd counting methods have been developed and tested on public datasets that are primarily outdoor scenes or general crowd environments, such as ShanghaiTech \cite{Zhu2020}, UCF-QNRF \cite{ucf_qnrf_dataset}, and NWPU-Crowd \cite{wang2020nwpu}. These datasets are not representative enough of the visual properties and problems of the interior of a bus. Another drawback of many studies is the estimation of crowd density instead of the identification of the individuals who make up the crowd in video frames, which can cause double counting if used in a video-based application and may lead to unstable counts \cite{hassan2023crowd}.

\section{Motivation and Objectives}

Bus interior scenes have special characteristics compared to general crowd scenes. Occupants often obstruct one another in the frame, cameras are often placed on the ceiling creating the illusion of perspective, and passengers' heads often seem to be small in photographs \cite{tezcan2024fisheye}. Thus, it is important to test if head detection can accurately estimate the number of passengers in a bus environment, explore optimal head detection models to detect small heads, and evaluate if head detection combined with object tracking can enhance the stability of counting and avoid multiple counting of passengers among different frames \cite{occupancy2026mot}. It is also important to investigate density estimation-based counting as a method to handle counting in scenarios where individual heads are difficult to accurately localise. Based on these considerations, the specific objectives are to investigate whether detecting people's heads is a valid method for counting the number of people in the interior of a bus; to study different object detection models and select the best model to detect passenger heads in real bus scenes; to investigate the use of density estimation-based counting as a method to handle counting in scenarios where individual heads are difficult to accurately localise; to develop a density decision method that can differentiate between bus interior scenes that contain only a few people and those that have more people; and to develop an adaptive hybrid counting model that can integrate the results of head detection and density estimation to obtain an accurate count of passengers.